% Generated by analysis.extended_report from measured extended results; do not edit.
\providecommand{\ExtOneRawAggregationMin}{2.4}
\providecommand{\ExtOneRawAggregationMax}{55.4}
\providecommand{\ExtOneRawRegenerationMin}{-6.1}
\providecommand{\ExtOneRawRegenerationMax}{24.2}
\providecommand{\ExtOneRawTotalMin}{-1.7}
\providecommand{\ExtOneRawTotalMax}{79.6}
\providecommand{\ExtOneRawAggregationPositive}{5}
\providecommand{\ExtOneGzipAggregationMin}{3.3}
\providecommand{\ExtOneGzipAggregationMax}{65.4}
\providecommand{\ExtOneGzipRegenerationMin}{-6.6}
\providecommand{\ExtOneGzipRegenerationMax}{7.0}
\providecommand{\ExtOneGzipTotalMin}{-2.2}
\providecommand{\ExtOneGzipTotalMax}{72.5}
\providecommand{\ExtOneGzipAggregationPositive}{5}
\providecommand{\ExtTwoRawAggregationMin}{2.2}
\providecommand{\ExtTwoRawAggregationMax}{52.8}
\providecommand{\ExtTwoRawRegenerationMin}{-6.2}
\providecommand{\ExtTwoRawRegenerationMax}{23.7}
\providecommand{\ExtTwoRawTotalMin}{-2.5}
\providecommand{\ExtTwoRawTotalMax}{76.5}
\providecommand{\ExtTwoRawAggregationPositive}{5}
\providecommand{\ExtTwoGzipAggregationMin}{3.7}
\providecommand{\ExtTwoGzipAggregationMax}{66.4}
\providecommand{\ExtTwoGzipRegenerationMin}{-6.7}
\providecommand{\ExtTwoGzipRegenerationMax}{5.9}
\providecommand{\ExtTwoGzipTotalMin}{-2.3}
\providecommand{\ExtTwoGzipTotalMax}{72.4}
\providecommand{\ExtTwoGzipAggregationPositive}{5}
\providecommand{\ExtTiles}{45{,}689}
\providecommand{\ExtCheckedAddresses}{91{,}378}
\providecommand{\ExtRoundtripDifferences}{0}
\providecommand{\ExtAggregateBlocks}{9{,}042}

Vector tiles are independently encoded display units, so tiling sets the scope over
which an encoder can exploit redundancy. We study spatial encoding-unit granularity
on five heterogeneous datasets using MVT and two versions of columnar MapLibre Tiles.
Five variants separate exact aggregation of conventional tile feature multisets from
regeneration of larger tiles at matched logical zoom. One pinned Rust implementation
encodes both MLT versions; geographic block matching gives an additive decomposition
of byte changes. For $4\times4$ fixed-content aggregation, raw MLT v1 byte reductions
range from \ExtOneRawAggregationMin\% to \ExtOneRawAggregationMax\%, and v2 from
\ExtTwoRawAggregationMin\% to \ExtTwoRawAggregationMax\%, across datasets. With
independent per-tile gzip, the ranges are
\ExtOneGzipAggregationMin--\ExtOneGzipAggregationMax\% and
\ExtTwoGzipAggregationMin--\ExtTwoGzipAggregationMax\%, respectively. Regenerated
$4\times4$ units give total gzip reductions of
\ExtOneGzipTotalMin--\ExtOneGzipTotalMax\% for v1 and
\ExtTwoGzipTotalMin--\ExtTwoGzipTotalMax\% for v2; negative reductions denote increases.
Exact multiset checks, source-feature provenance and version-specific stream statistics
distinguish encoding-unit effects from changes in stored content. Format comparisons
are conditional on granularity and generation policy; they do not establish
client-side transfer or rendering benefits.
